\documentclass[aspectratio=169]{beamer}
\usepackage{iftex}
\ifPDFTeX\usepackage[T1]{fontenc}\usepackage{lmodern}\fi
\usepackage[czech]{babel}\usepackage{listings}
\newcommand{\CVUTTemplatePath}{../../../utils/presentation-template}
\usepackage{\CVUTTemplatePath/beamerthemeCVUT}
\setbeamertemplate{footline}[frame number]\beamertemplatenavigationsymbolsempty
\newcommand{\LANG}{CZ}\newcommand{\FRONTPAGE}{dark}
\title{Sdílené proměnné a komunikace mezi částmi programu}
\subtitle{Základy programování\\Matouš Cejnek\\U12110, FS, ČVUT v Praze}
\author{}\institute{}\date{}
\renewcommand{\headlineA}{}\renewcommand{\headlineB}{}\renewcommand{\headlineC}{}
\definecolor{CodeBlue}{RGB}{38,83,125}\definecolor{SoftGray}{RGB}{242,244,247}
\definecolor{CodeGreen}{RGB}{38,125,94}\definecolor{WarningRed}{RGB}{170,45,45}
\setbeamerfont{normal text}{size=\large}\setbeamerfont{frametitle}{size=\LARGE}
\setbeamercolor{block title}{fg=white,bg=CodeBlue}\setbeamercolor{block body}{fg=black,bg=SoftGray}
\newcommand{\term}[1]{\textcolor{CodeBlue}{\textbf{#1}}}
\lstdefinestyle{python}{language=Python,basicstyle=\ttfamily\normalsize,
 keywordstyle=\color{CodeBlue}\bfseries,stringstyle=\color{CodeGreen},
 commentstyle=\color{gray},showstringspaces=false,columns=fullflexible,keepspaces=true}
\begin{document}
\begin{frame}\clearpage\thispagestyle{empty}\titlepage\end{frame}

\begin{frame}{Co si dnes odnesete?}
\begin{itemize}
\item Rozlišit lokální, globální a sdílený stav.
\item Předávat data explicitně mezi funkcemi.
\item Předvídat změny sdílených mutabilních objektů.
\item Vyhnout se skrytým závislostem a mutabilním výchozím hodnotám.
\item Zapouzdřit související stav a chování do objektu.
\end{itemize}
\end{frame}

\begin{frame}{Co je stav programu?}
\begin{itemize}
\item Hodnoty, které program v daném okamžiku uchovává.
\item Lokální proměnné právě volaných funkcí.
\item Obsah objektů, kolekcí, souborů nebo databáze.
\item Stav ovlivňuje výsledek dalších operací.
\end{itemize}
\begin{block}{Příklad}Zůstatek účtu před výběrem je součástí stavu.\end{block}
\end{frame}

\begin{frame}{Komunikace mezi částmi programu}
\begin{itemize}
\item Parametry nesou vstupy do funkce.
\item Návratová hodnota nese výsledek ven.
\item Objekt může nést související stav mezi voláními metod.
\item Soubor, fronta nebo databáze propojuje vzdálenější části.
\item Globální proměnná je pouze jedna z možností --- často problematická.
\end{itemize}
\end{frame}

\begin{frame}[fragile]{Lokální proměnná patří jednomu volání}
\begin{lstlisting}[style=python]
def calculate_price(amount, unit_price):
    subtotal = amount * unit_price
    return subtotal

price = calculate_price(3, 120)
\end{lstlisting}
\begin{itemize}\item \texttt{subtotal} po návratu funkce není zvenku dostupná.\end{itemize}
\end{frame}

\begin{frame}{Rozsah platnosti jména}
\begin{description}
\item[Local] Jména uvnitř aktuální funkce.
\item[Enclosing] Jména v obalujících funkcích.
\item[Global] Jména na úrovni modulu.
\item[Built-in] Vestavěná jména, například \texttt{len}.
\end{description}
\begin{block}{LEGB}Python hledá jméno v tomto pořadí.\end{block}
\end{frame}

\begin{frame}[fragile]{Čtení globální hodnoty}
\begin{lstlisting}[style=python]
VAT_RATE = 0.21

def price_with_vat(price):
    return price * (1 + VAT_RATE)
\end{lstlisting}
\begin{itemize}
\item Konstantní konfigurace může být na úrovni modulu čitelná.
\item Funkce však má skrytou závislost na \texttt{VAT\_RATE}.
\end{itemize}
\end{frame}

\begin{frame}[fragile]{Zápis do globální proměnné}
\begin{lstlisting}[style=python]
score = 0

def add_point():
    global score
    score += 1
\end{lstlisting}
\begin{itemize}
\item \texttt{global} říká, že zápis míří do jména modulu.
\item Výsledek funkce závisí na historii předchozích volání.
\end{itemize}
\end{frame}

\begin{frame}{Proč měnitelný globální stav bolí?}
\begin{itemize}
\item Závislost není vidět v parametrech funkce.
\item Na pořadí volání záleží.
\item Testy si mohou navzájem měnit prostředí.
\item Změnu může provést mnoho míst programu.
\item Chyba se projeví daleko od svého původu.
\end{itemize}
\end{frame}

\begin{frame}[fragile]{Explicitní předání stavu}
\begin{lstlisting}[style=python]
def add_point(score):
    return score + 1

score = 0
score = add_point(score)
\end{lstlisting}
\begin{block}{Výhoda}Vstup i výstup jsou vidět v rozhraní funkce.\end{block}
\end{frame}

\begin{frame}{Čistá funkce}
\begin{itemize}
\item Výsledek závisí pouze na argumentech.
\item Nemění stav mimo své lokální prostředí.
\item Pro stejné vstupy vrací stejný výsledek.
\item Snadno se testuje a skládá s dalšími funkcemi.
\item Ne každá užitečná funkce může být čistá, ale čisté jádro pomáhá.
\end{itemize}
\end{frame}

\begin{frame}{Vedlejší účinek}
\begin{itemize}
\item Změna objektu předaného argumentem.
\item Zápis do souboru nebo databáze.
\item Výpis na obrazovku či komunikace po síti.
\item Změna globálního stavu.
\end{itemize}
\begin{block}{Není automaticky špatně}Vedlejší účinek má být záměrný, pojmenovaný a oddělený.\end{block}
\end{frame}

\begin{frame}{Jméno odkazuje na objekt}
\begin{itemize}
\item Při přiřazení se objekt obvykle nekopíruje.
\item Dvě jména mohou odkazovat na tentýž objekt.
\item Změna mutabilního objektu je viditelná přes všechny odkazy.
\item Nové přiřazení jména není totéž jako změna objektu.
\end{itemize}
\end{frame}

\begin{frame}[fragile]{Dvě jména, jeden seznam}
\begin{lstlisting}[style=python]
first = [1, 2]
second = first
second.append(3)

print(first)   # [1, 2, 3]
\end{lstlisting}
\begin{itemize}\item \texttt{first is second} má hodnotu \texttt{True}.\end{itemize}
\end{frame}

\begin{frame}{Mutabilní a imutabilní objekty}
\begin{description}
\item[Mutabilní] Seznam, slovník, množina a většina vlastních instancí.
\item[Imutabilní] Číslo, řetězec, \texttt{tuple} z imutabilních prvků.
\end{description}
\begin{itemize}
\item U imutabilní hodnoty operace vytvoří nový objekt.
\item U mutabilní hodnoty může operace změnit existující objekt.
\end{itemize}
\end{frame}

\begin{frame}[fragile]{Argument se předává jako odkaz na objekt}
\begin{lstlisting}[style=python]
def add_item(items, item):
    items.append(item)

values = [1, 2]
add_item(values, 3)
print(values)  # [1, 2, 3]
\end{lstlisting}
\begin{itemize}\item Funkce i volající vidí tentýž seznam.\end{itemize}
\end{frame}

\begin{frame}[fragile]{Nové přiřazení argumentu objekt nezmění}
\begin{lstlisting}[style=python]
def replace(items):
    items = [0]

values = [1, 2]
replace(values)
print(values)  # [1, 2]
\end{lstlisting}
\begin{itemize}\item Změnilo se pouze lokální jméno \texttt{items}.\end{itemize}
\end{frame}

\begin{frame}[fragile]{Past: mutabilní výchozí argument}
\begin{lstlisting}[style=python]
def remember(value, history=[]):
    history.append(value)
    return history

remember("A")  # ["A"]
remember("B")  # ["A", "B"]
\end{lstlisting}
\begin{itemize}\item Výchozí objekt vznikne jednou při definici funkce.\end{itemize}
\end{frame}

\begin{frame}[fragile]{Bezpečný vzor s \texttt{None}}
\begin{lstlisting}[style=python]
def remember(value, history=None):
    if history is None:
        history = []
    history.append(value)
    return history
\end{lstlisting}
\begin{itemize}\item Každé volání bez argumentu vytvoří nový seznam.\end{itemize}
\end{frame}

\begin{frame}{Kopírování není vždy jednoduché}
\begin{itemize}
\item Mělká kopie vytvoří novou vnější kolekci.
\item Vnořené objekty mohou zůstat sdílené.
\item Hluboká kopie kopíruje i vnořenou strukturu, ale je dražší.
\item Lepší než kopírovat vše je určit jasného vlastníka dat.
\end{itemize}
\end{frame}

\begin{frame}{Jasné vlastnictví stavu}
\begin{itemize}
\item Jedna část programu odpovídá za změny konkrétních dat.
\item Ostatní části žádají o operaci přes definované rozhraní.
\item Pravidla změny jsou soustředěna na jednom místě.
\item Stav se lépe validuje, testuje a ladí.
\end{itemize}
\end{frame}

\begin{frame}[fragile]{Objekt jako nosič stavu}
\begin{lstlisting}[style=python]
class ScoreTracker:
    def __init__(self):
        self.score = 0

    def add_point(self):
        self.score += 1
\end{lstlisting}
\begin{itemize}\item Stav je stále měnitelný, ale má konkrétního vlastníka.\end{itemize}
\end{frame}

\begin{frame}{Oddělme výpočet od vstupu a výstupu}
\begin{itemize}
\item Čisté funkce provádějí transformace dat.
\item Okraj programu čte vstupy a zapisuje výsledky.
\item Stavový objekt koordinuje související operace.
\item Menší počet míst s vedlejšími účinky zjednodušuje testy.
\end{itemize}
\end{frame}

\begin{frame}{Sdílený stav a souběžnost}
\begin{itemize}
\item Dvě činnosti mohou stav číst a měnit současně.
\item Výsledek pak může záviset na pořadí kroků.
\item Synchronizace řeší přístup, ale přidává složitost.
\item I v jednovláknovém programu pomáhá sdílení omezovat.
\end{itemize}
\begin{block}{Nad rámec tématu}Vlákna a asynchronní programování mají další pravidla.\end{block}
\end{frame}

\begin{frame}{Rozhodovací vodítko}
\begin{enumerate}
\item Stačí hodnota pro jeden výpočet? Použijte parametr a návrat.
\item Patří stav a operace k sobě? Zvažte objekt.
\item Musí data přežít běh programu? Použijte soubor či databázi.
\item Je globální stav opravdu nutný? Omezte zápis a jasně jej pojmenujte.
\end{enumerate}
\end{frame}

\begin{frame}{Modul je také hranice}
\begin{itemize}
\item Modul může skrýt pomocná jména a nabídnout veřejné funkce.
\item Import modulu nevytváří novou kopii jeho globálního stavu.
\item Změna modulu je viditelná všem jeho uživatelům v procesu.
\item Veřejné rozhraní má být menší než vnitřní implementace.
\end{itemize}
\end{frame}

\begin{frame}{Více hodnot předáme jako celek}
\begin{itemize}
\item Několik souvisejících výsledků může vrátit \texttt{tuple}.
\item Slovník pojmenuje položky, ale dovoluje snadné překlepy.
\item Datová třída dává hodnotám názvy a jasný společný typ.
\item Volba struktury vyjadřuje význam předávaných dat.
\end{itemize}
\end{frame}

\begin{frame}{Rovnost a totožnost}
\begin{description}
\item[\texttt{==}] Mají objekty stejnou hodnotu?
\item[\texttt{is}] Jde o tentýž objekt?
\end{description}
\begin{itemize}
\item Dva různé seznamy mohou být hodnotově stejné.
\item \texttt{is} používáme typicky pro porovnání s \texttt{None}.
\end{itemize}
\end{frame}

\begin{frame}{Imutabilní návrh omezuje překvapení}
\begin{itemize}
\item Místo změny existující hodnoty vytvoříme hodnotu novou.
\item Původní stav zůstane dostupný ostatním částem programu.
\item Tok změn je vidět v přiřazení návratové hodnoty.
\item Cena může být vyšší spotřeba paměti nebo kopírování.
\end{itemize}
\end{frame}

\begin{frame}{Závislost lze předat parametrem}
\begin{itemize}
\item Funkce nemusí sama hledat globální konfiguraci nebo službu.
\item Volající předá potřebnou hodnotu či spolupracující objekt.
\item V testu lze předat jednoduchou kontrolovanou náhradu.
\item Rozhraní jasně ukazuje, co funkce potřebuje ke své práci.
\end{itemize}
\end{frame}

\begin{frame}{Kdo smí stav měnit?}
\begin{itemize}
\item Čtení lze nabídnout široce, zápis omezit na několik operací.
\item Validace proběhne před každou povolenou změnou.
\item Metoda pojmenovává záměr lépe než přímý zápis atributu.
\item Audit změn je snazší, když vedou přes jedno rozhraní.
\end{itemize}
\end{frame}

\begin{frame}{Shrnutí}
\begin{itemize}
\item Explicitní tok dat usnadňuje porozumění i testování.
\item Globální měnitelný stav vytváří skryté závislosti.
\item Sdílený mutabilní objekt se mění přes všechny své odkazy.
\item Výchozí mutabilní hodnoty nahrazujeme vzorem s \texttt{None}.
\item Objekt může být jasným vlastníkem souvisejícího stavu.
\end{itemize}
\end{frame}
\end{document}
